\chapter{Speculative decoding and attention kernels}
\label{evod38:chapter}

Suppose the target next-token law is $(.55,.25,.15,.05)$ while a cheaper
draft proposes from $(.45,.35,.15,.05)$. Accepting every proposal is fast
but wrong: it overproduces token one and underproduces token zero. Rejecting
some proposals and sampling from the original target distribution after
rejection also gives the wrong mixture. The correction must supply precisely
the target mass left uncovered by accepted proposals.

This chapter finishes two distinct engine mechanisms. Speculation changes
which future computations are attempted before their need is known. Online
attention changes how the exact normalization/value sum is accumulated.
Both retain the actual Chapter 23/35 Hermon model and Chapter 36 page
lifetime contracts. Our complete-stack implementations are CPU references;
the title's kernel discussion supplies derivations and backend obligations,
not an unmeasured GPU speed claim.

\section{Construct a target-preserving one-token sampler}
\index{Speculative decoding}\index{Modified rejection sampling}\index{Residual correction}\index{Total variation}
\label{evod38:sampling}
Fix a finite vocabulary and prefix $h$. Let $p(x\mid h)$ be the intended
target distribution and $q(x\mid h)$ the actual draft proposal law, both
nonnegative and normalized. Draw $X\sim q$ and independent $U\sim
\mathrm{Uniform}[0,1)$. Accept $X$ when
\[
 U<\alpha(X),\qquad \alpha(x)=\min(1,p_x/q_x)\quad(q_x>0).
\]
A $q_x=0$ token is never proposed, so no division is needed there.
Define $d_x=(p_x-q_x)_+$ and $R=\sum_x d_x$. If a proposal is rejected,
draw a fresh token from $r_x=d_x/R$. When $R=0$, $p=q$ and rejection has
zero probability in ideal arithmetic; the reference returns $p$ as an
unused correction fallback.

\begin{proposition}[One-position exact law]
Under those distributions and independent exact random draws, the emitted
token $Y$ has distribution $p$.
\end{proposition}
\begin{proof}
The joint probability of proposing and accepting token $x$ is
$q_x\alpha(x)=\min(p_x,q_x)$. Since $p$ and $q$ sum to one, their total
positive and negative difference masses agree. Thus
\[
 1-\sum_x\min(p_x,q_x)=\sum_x(p_x-q_x)_+=R.
\]
The probability of rejecting any proposal and then emitting $x$ is
$Rr_x=d_x$. These disjoint paths sum to
$\min(p_x,q_x)+(p_x-q_x)_+=p_x$.
The $R=0$ case emits only accepted proposals, which already have law $p$.
\end{proof}

This is the modified rejection mechanism used in speculative decoding
\cite{evod38:leviathan,evod38:chen}. The reference uses $p$ for target and
$q$ for draft; some papers reverse those letters. Track roles, not letters.

\begin{center}\input{\Generated/acceptance.tex}\end{center}
\Plate{01-mass}{Accepted proposal mass totals .90. The missing .10 belongs entirely to token zero, so every rejection correction selects that token in this example.}{evod38:mass}

Here $R=.10$ and $r=(1,0,0,0)$. If we incorrectly used $p$ itself after
rejection, token zero would receive $.45+.10(.55)=.505$, not .55.
The negative test makes this incorrect fallback fail visibly. Support
gaps are allowed: if $p=(.5,.5)$ and $q=(1,0)$, acceptance supplies .5 at
token zero and correction supplies .5 at token one.

The expected acceptance is
\[
 A=\sum_x\min(p_x,q_x)
 =1-\tfrac12\sum_x|p_x-q_x|=1-\mathrm{TV}(p,q).
\]
This measures distribution overlap at a specified prefix. It does not
measure draft cost, cache allocation, communication or latency.

\section{A block contains conditional proposals, not independent rows}
\index{Autoregressive distribution}\index{Conditional proposal}\index{Sampling!random streams}
\label{evod38:blocks}
The draft proposes $x_1,\ldots,x_k$ sequentially, retaining each actual
law $q_i(\cdot\mid h,x_{<i})$. The target computes
$p_i(\cdot\mid h,x_{<i})$ for the same candidate prefixes. Chapter 35's
offset-masked causal chunk allows these target rows to be evaluated
together. The verification draws use independent uniforms; at the first
rejection $i$, emit accepted $x_{<i}$ followed by a correction $y_i$, and
stop. Candidate $x_{i+1}$ was conditioned on $x_i$, not $y_i$.

\Plate{02-dependencies}{Only a contiguous accepted prefix is reusable. A rejection invalidates all descendants conditioned on the rejected proposal, even when their tensors are already computed.}{evod38:dependencies}

\Listing{verify}

\begin{proposition}[Autoregressive law across blocks]
Suppose target/draft callbacks depend only on the declared token prefix,
both use the same token identities, and each proposal/verification draw is
fresh with the prescribed conditional law. Repeated corrected blocks,
truncated to a requested output length, have target joint distribution
$\prod_j p(y_j\mid h,y_{<j})$.
\end{proposition}
\begin{proof}
At the first uncommitted position, the one-position construction supplies
the target conditional. If accepted, the next candidate was proposed at
exactly the new committed prefix, so the same argument applies to that
position. If rejected, the correction has completed the one-position
mixture; later candidates are discarded and a fresh block starts at the
corrected committed prefix. Thus every transition supplies the target law
at its actual prefix. Multiplying conditional transition probabilities
proves the joint law by induction over emitted positions.
\end{proof}

Merely passing a list of unconditional $p_i,q_i$ arrays is not a general
autoregressive implementation. The supplied package's list-based example
is retained as the one-position algebraic lead and extended into explicit
prefix callbacks and recorded proposal laws. Tests enumerate a two-token,
two-symbol history-dependent case exactly, then check a seeded 20,000-run
joint-frequency experiment with a predeclared .015 absolute tolerance.
That finite statistical check supports the implementation; the proof, not
one histogram, establishes the ideal law.

\subsection{Identity, transforms and randomness}
The target $p$ must include the intended temperature, grammar mask, top-k
or nucleus transformation and renormalization. The ratio must use the law
from which the draft actually sampled. The draft may use a \emph{different}
temperature or support restriction: the proof does not require $p=q$ or
identical transforms, only correctly recorded laws on compatible token IDs.
A tokenizer mismatch or probabilities from an earlier prefix invalidate it.

Exact distribution equality does not mean token-for-token equality under
the same seed as target-only sampling: the algorithms consume different
random draws. Record RNG implementation, seed/stream policy, target/draft
artifacts, transforms and numerical representation for reproducibility.
Floating-point normalization and pseudorandom sampling approximate the
ideal proof. The reference rejects nonfinite/negative vectors, mismatched
vocabulary sizes and impossible zero-mass proposals; it rescales finite
unnormalized inputs before summing to avoid overflow.

\section{Verification state is tentative ownership}
\index{Speculative decoding!rollback}\index{KV cache!tentative fork}
\label{evod38:commit}
The native adapter starts from a nonempty consumed prefix and obtains its
initial next-token logits. It forks the Chapter 36 live handle into a
same-tenant verification owner and consumes the draft block there. If the
prefix has length $t$, the initial logits supply $p_1$; the row computed
after consuming $x_i$ supplies $p_{i+1}$. Shifting these rows by one would
verify each proposal against the wrong conditional.

After sampling the accepted/corrected output, it releases the entire
tentative fork and recomputes only the committed tokens on the original
handle. This conservative strategy avoids shortening shared physical pages.
It is deliberately expensive, but a simple correctness target for a future
transactional suffix adoption operation.

\Plate{03-rollback}{The live table survives verification unchanged. Releasing tentative ownership reclaims only unreferenced pages; committed recomputation then advances the live cache.}{evod38:rollback}
\Listing{native_speculative}

In this adapter, every returned token has also been consumed into KV.
Callers must not consume it again. This differs from Chapter 37's final
unconsumed-output convention and requires an explicit integration adapter.
The no-bonus variant returns at most $k$ tokens. Standard variants can
emit a target bonus after all $k$ proposals are accepted, provided their
distribution comes from the full accepted prefix and its cache accounting
is correct. Omitting a bonus changes throughput, not the corrected target law.

On capacity failure, the verification fork is released in a \texttt{finally}
block and the live prefix remains unchanged. Tests force OOM on a shared
partial tail and compare payload seals, refcounts and live-handle count.
On successful commits, tests reconstruct the whole prefix with the original
full model and compare every layer's K/V. Local serial cleanup does not
certify concurrent cancellation or crash recovery.

\section{Acceptance is not a speed measurement}
\index{Acceptance probability}\index{Bonus token}\index{Tail-sum identity}
\label{evod38:cost}
Let $N$ be emitted tokens per no-bonus iteration. It emits at least one;
for $j\ge2$, $N\ge j$ exactly when the first $j-1$ proposals are accepted.
The tail-sum identity gives
\[
 \mathbb E[N]=\sum_{j=1}^k\Pr(N\ge j)
 =1+\sum_{i=1}^{k-1}\Pr(\text{first }i\text{ accepted}).
\]
Under the additional approximation that acceptance events are independent
with constant probability $A$, this becomes
$\sum_{i=0}^{k-1}A^i=(1-A^k)/(1-A)$ for $A\ne1$, and $k$ for $A=1$.
History-dependent real acceptance generally violates that approximation.

For $k=3,A=.9$, the no-bonus model emits $1+.9+.81=2.71$ tokens per
iteration on average. A bonus variant adds $.9^3=.729$, giving 3.439.
With draft cost $kc_d$, target verification cost $c_v(k)$ and bookkeeping
$c_b$, estimated time per emitted token is
\[
 \frac{kc_d+c_v(k)+c_b}{\mathbb E[N]}.
\]
Compare this with measured target-only per-token time. Do not set
$c_v(k)$ equal to a scalar decode cost without profiling. Our adapter also
does prefix replay and committed recomputation; it has no speedup claim.

\section{Attention can avoid a global score buffer}
\index{Online softmax}\index{Attention!tiled reduction}\index{Masked attention!empty row}
\label{evod38:attention}
For one query $q\in\mathbb R^{d}$, keys $k_j\in\mathbb R^d$, values
$v_j\in\mathbb R^{d_v}$ and allowed index set $J$,
\[
 s_j=q^\top k_j/\sqrt d,\qquad
 y=\frac{\sum_{j\in J} e^{s_j}v_j}{\sum_{j\in J}e^{s_j}}.
\]
Direct exponentiation overflows for large scores. Subtracting a common
maximum preserves the ratio. To traverse tiles, retain summary
\[
 m_J=\max_{j\in J}s_j,\quad
 \ell_J=\sum_{j\in J}e^{s_j-m_J},\quad
 a_J=\sum_{j\in J}e^{s_j-m_J}v_j.
\]
Here $m$ and $\ell$ are scalars, $a\in\mathbb R^{d_v}$. For several
heads/queries, maintain one such state per row. Values need not be positive;
the normalizer is positive for every nonempty allowed set.

For disjoint tiles A and B, use $m=\max(m_A,m_B)$ and
\begin{align*}
 \ell&=e^{m_A-m}\ell_A+e^{m_B-m}\ell_B,\\
 a&=e^{m_A-m}a_A+e^{m_B-m}a_B.
\end{align*}
This rescales the \emph{old accumulator too}. Updating only the normalizer
when a larger maximum arrives produces a wrong weighted value sum.

\begin{proposition}[Online-softmax invariant]
After any collection of disjoint allowed tiles, the merged summary equals
$(m_J,\ell_J,a_J)$ for their union in real arithmetic. Consequently
$a_J/\ell_J$ is the dense attention result. Merge is associative for
nonempty sets under these exact summaries.
\end{proposition}
\begin{proof}
Each old term obeys
$e^{m_A-m}e^{s_j-m_A}=e^{s_j-m}$, and similarly for B. Summing the
rescaled normalizers and vector accumulators therefore gives exactly the
definitions for the union. The maximum also equals the union maximum.
The union operation is associative, so any merge tree gives the same
summary. Division cancels the common $e^{-m}$ factor. Induction over tiles
proves the online recurrence.
\end{proof}

\Plate{04-online}{A larger second-tile maximum rescales the first tile by one half. The merged output is six, exactly the dense weighted mean.}{evod38:onlineplate}

The empty-set identity needs care: represent it by $m=-\infty,\ell=0,a=0$
but never evaluate $-\infty-(-\infty)$. The implementation substitutes a
safe finite shift for an entirely masked row; all its exponential weights
are zero. It returns zero attention when the full allowed set is empty,
matching the earlier chapters' explicit policy. For a nonempty row at least
one shifted weight is one, so the denominator cannot vanish ideally.
Finite inputs that overflow the score dot product are rejected separately.

\Listing{online_attention}

Floating-point merge is not bitwise associative. Tile sizes and reduction
order can change low-order bits. Tests compare tile sizes 1, 2, 3, 7 and
oversized tiles against an independent dense implementation and PyTorch
2.10 SDPA, with explicit boolean masks and inference dropout zero
\cite{evod38:sdpa}. They include extreme finite scores, initially masked
tiles and all-masked rows. A normalization proof does not by itself validate
every device backend.

\section{Preserve offset causality and the complete model}
\index{Offset causal mask}\index{Attention!paged traversal}
\label{evod38:native}
With $t$ past tokens and a $c$-token chunk, local query $i$ is absolute
position $t+i$, so key $j$ is allowed exactly when $j\le t+i$. A physical
page boundary does not reset that ordinal. For $t=3,c=2$, the allowed row
counts are four and five, not one and two from an upper-left triangle.

\Plate{05-mask}{The same causal mask applies across tiles and scattered pages. The blocked future cell belongs to the new chunk, not to an arbitrary physical page boundary.}{evod38:mask}

\texttt{online\_consume} replaces only the attention reduction in the
actual native stack: learned absolute embeddings, normalized QKV,
multi-head output projection, residuals, FFN, final normalization and LM
head remain. It preserves all layer caches and the token/owner/artifact
contract. Whole-function extraction keeps printed code bound to the tested
implementation.

\Listing{online_consume}
\begin{center}\input{\Generated/attention.tex}\end{center}

The table compares continuation after a three-token prefix across tile
sizes against the original full causal model. Additional tests vary every
split of the six-token record and compare every layer's K/V. The separate
\texttt{paged\_attention} function visits Chapter 36 pages in table order
without flattening them for the attention reduction. Its queries and masks
use logical ordinals; a scattered physical layout is a test fixture.
That page visitor is not yet the complete model's in-place paged append
kernel. The distinction is an explicit remaining backend integration task.

\section{IO awareness, fusion and evidence}
\index{FlashAttention}\index{Kernel fusion}\index{Memory traffic!model}
\label{evod38:io}
FlashAttention's tiled exact-attention algorithm and on-chip fusion provide
the primary systems comparison~\cite{evod38:flash}; FlashAttention-2 further
changes parallelism/work partitioning~\cite{evod38:flash2}. These are backend
designs, not a claim that any Python tiled loop has the same performance.
Fusion can reduce global intermediate transfers while increasing register
pressure, reducing occupancy or restricting shapes. Measure the combined
page size, head geometry, dtype, masking and batch distribution.

For a deliberately narrow one-query traffic model, assume each Q/K/V input
is read once and output written once, with element width $b$ bytes. The
base payload is
\[
 b(d+Td+Td_v+d_v).
\]
An explicitly materialized score vector that is written once and read once
adds $2Tb$ bytes. We exclude probability buffers, cache effects, repeated
tile loads, metadata, weights and workspaces. This model cannot stand in
for FlashAttention's multi-query IO complexity theorem or measured HBM traffic.
\begin{center}\input{\Generated/traffic.tex}\end{center}
For $T=4096,d=d_v=128,b=2$, the isolated saving is 16,384 bytes,
less than one percent of this one-query payload model. It would be
misleading to turn that number into a universal speedup. Multi-query
score matrices and their reuse require a different calculation.

Low-precision tiled attention must satisfy Chapter 37's numerical/task
gates; exact algebra does not erase storage or reduction error. Speculative
verification must use the intended target policy, including its actual
precision. A quantized target's exact sampling law is not automatically
the original float64 model's law.

A credible kernel study proceeds from semantics and adversarial parity to
declared traffic/work models, then synchronized hardware profiles and
finally service measurements. Record warmup, hardware/software versions,
repetitions, shape distribution, offered load, TTFT/ITL tails, throughput,
memory high-water and rejected/preempted work. Correctness, model arithmetic,
CPU experiments, device measurements and end-to-end SLO results are distinct
evidence levels.

\section{Selective glossary}
\begin{description}
\item[Proposal law] Actual conditional distribution used by the draft to draw a candidate.
\item[Residual correction] Normalized positive target-minus-draft mass used after rejection.
\item[Accepted prefix] Contiguous candidates verified before the first rejection.
\item[Tentative KV] Computed state owned by verification, not yet part of committed history.
\item[Online summary] Maximum, shifted exponential sum and shifted value accumulator for processed allowed keys.
\item[Exact attention] The same mathematical softmax/value map, with finite-precision differences governed by a numerical contract.
\item[IO model] Declared transfer-count assumptions; it is distinct from observed hardware traffic or service latency.
\end{description}

\section{Twenty worked exercises}
\label{evod38:exercises}
\begin{enumerate}
\item \textbf{Equal laws.} $p=q$? \emph{Solution:} acceptance one;
correction branch has zero probability.
\item \textbf{Support gap.} $p=(.5,.5),q=(1,0)$?
\emph{Solution:} token zero accepted with .5; rejection corrects to token one.
\item \textbf{Compute overlap.} Supplied four-token example?
\emph{Solution:} $.45+.25+.15+.05=.90=1-\mathrm{TV}$.
\item \textbf{Reject the wrong fallback.} Sample $p$ after rejection?
\emph{Solution:} token-zero probability becomes .505 instead of .55.
\item \textbf{Conditional rows.} Which prefix supplies $p_3$?
\emph{Solution:} original $h$ followed by proposed $x_1,x_2$.
\item \textbf{First rejection.} Why discard $x_3$ when $x_2$ rejects?
\emph{Solution:} $x_3$ used a history that was not committed.
\item \textbf{Transforms.} Must draft temperature equal target temperature?
\emph{Solution:} no; ratio must use actual $q$ and intended transformed $p$.
\item \textbf{Reproduce seeds.} Does same seed force target-only equality?
\emph{Solution:} no; draw counts/order differ even when laws agree.
\item \textbf{Shift logits.} Which row verifies first proposal?
\emph{Solution:} last old-prefix logits, before consuming that proposal.
\item \textbf{Commit state.} Can a rejected suffix stay in live KV?
\emph{Solution:} no; discard tentative ownership and install only committed history.
\item \textbf{Compare API clocks.} Returned native speculative tokens?
\emph{Solution:} already consumed; avoid consuming them a second time.
\item \textbf{Expected count.} No bonus, $k=3,A=.9$ independent?
\emph{Solution:} $1+.9+.81=2.71$ emitted, not $3(.9)$.
\item \textbf{Add a bonus.} Same approximation? \emph{Solution:}
add $.729$ to get 3.439; this requires the extra target conditional.
\item \textbf{Merge summaries.} Plate's $(m,\ell,a)$?
\emph{Solution:} $(\log4,1.75,10.5)$; ratio six.
\item \textbf{Find an unstable update.} Rescale $\ell$ but not $a$?
\emph{Solution:} plate would use $a=13$ instead of 10.5, giving a wrong mean.
\item \textbf{Empty masks.} How avoid $-\infty-(-\infty)$?
\emph{Solution:} safe finite shift for empty rows, zero weights and explicit zero output.
\item \textbf{Offset mask.} Past three, chunk two?
\emph{Solution:} allowed counts four and five across all page tiles.
\item \textbf{Traffic.} $2Tb$ for $T=4096,b=2$?
\emph{Solution:} 16,384 isolated score-buffer bytes, not measured speed.
\item \textbf{Design a distribution test.} \emph{Rubric:} include zero
support, $p=q$, a history-dependent exact enumeration, seeded statistics
with fixed tolerances and a known wrong correction control.
\item \textbf{Design a backend release gate.} \emph{Rubric:} whole-stack
and per-layer parity, page-order/mask adversaries, precision/task gates,
rollback/completion fences, synchronized device profiles and service load curves.
\end{enumerate}

\section{Closing the two engine directions}
Chapters 30--34 supplied runtime/state-native references for a restricted
DOGMA teaching recurrence; Chapters 35--38 now supply causal Hermon
Transformer cache, allocation, scheduling and exact-attention references.
Neither establishes equal learned capability or universal performance
superiority. Evolutor's next layer must route typed work above those distinct
semantics without exposing either engine's mutable internals as a universal
state format. Chapter 39 owns that heterogeneous runtime boundary.
